% SPDX-License-Identifier: AGPL-3.0-or-later
% reflowtex capture, read before the document.
%
% The pipeline compiles a driver file (src/pipeline/driver.ts) that reads
% this file, then the document itself, unchanged: TeX chooses the class
% (behind a conditional or in a file the document inputs, as it may be),
% follows the inputs and expands the conditionals. Nothing here loads a
% package or a class: LaTeX wants \DocumentMetadata before any package, and
% the document's own options for the packages it loads must be the first
% ones. What has to wait for the class or for a package is installed by
% LaTeX's hooks: at the end of the class (floats and captions), when
% graphics is loaded (the graphics record), and at \begin{document}
% (reflowtex-capture-late.tex: everything that must see the preamble's
% definitions).

% ── The document's files, by names relative to its directory ────────────────
% TeX runs in a build directory of its own and finds the document's files
% through TEXINPUTS, which the build sets to the document's directories. A
% name that starts with ./ or ../ (\input{./sections/intro},
% \includegraphics{./figure}) TeX takes as it stands, relative to the build
% directory, and never looks for along TEXINPUTS. LaTeX looks for a name it
% did not find in each directory of \input@path in turn, whatever the name –
% \input, \include, \InputIfFileExists, a class or package, and graphicx's
% files unless the document sets \graphicspath – so TEXINPUTS's directories
% go there too, and a relative name resolves against the document's
% directory, as in a LaTeX run in that directory.
%
% What TeX writes goes to the build directory, which stands in for the
% document's: a name with a directory (\include{chapters/one} writes
% chapters/one.aux) needs that directory there, as the document's directory
% has it. So when TeX opens a file for writing in a directory the build
% directory lacks but one of the document's directories has, it is made
% first (directories only; never above the build directory). A directory the
% document's own lacks stays missing, and TeX says it cannot write there, as
% in a LaTeX run in the document's directory. (Backslashes and braces are
% built with string.char so TeX does not read them; the directories are
% printed as characters, whatever they hold.)
\directlua{
  local bs, open, close = string.char(92), string.char(123), string.char(125)
  local sep = os.type == "windows" and ";" or ":"
  local dirs = {}
  for d in string.gmatch(os.getenv("TEXINPUTS") or "", "[^" .. sep .. "]+") do
    d = d:gsub("/+$", "")
    if not (d == "" or d == ".") and lfs.isdir(d) then table.insert(dirs, d .. "/") end
  end
  if next(dirs) then
    tex.sprint(bs .. "expandafter" .. bs .. "def" .. bs .. "csname input@path" .. bs .. "endcsname" .. open)
    for _, d in ipairs(dirs) do tex.sprint(open); tex.sprint(-2, d); tex.sprint(close) end
    tex.sprint(close)
  end
  luatexbase.add_to_callback("find_write_file", function(_, name)
    local sub = string.match(name, "^(.+)/[^/]+$")
    if sub and not lfs.isdir(sub) then
      local parts, up = {}, false
      for part in string.gmatch(sub, "[^/]+") do table.insert(parts, part); up = up or part == ".." end
      for _, d in ipairs(up and {} or dirs) do
        if lfs.isdir(d .. sub) then
          local path = ""
          for _, part in ipairs(parts) do path = path .. part; lfs.mkdir(path); path = path .. "/" end
          break
        end
      end
    end
    return name
  end, "reflowtex: the document's directories in the build directory")
}

% ── No terminal: the run stays in nonstopmode ───────────────────────────────
% The pipeline runs TeX with -interaction=nonstopmode and stdin closed
% (src/pipeline/lualatex.ts): nobody answers TeX. A document that switches to
% a mode that asks – \errorstopmode, \scrollmode, as a class does after a
% quiet \batchmode check and \tracingall does – would have TeX read the
% terminal at its next error, find the end of the file and stop. The switch
% keeps nonstopmode instead, so the error is one TeX carries on past, as it
% is under nonstopmode; \batchmode is left as it is (it only writes less to
% the terminal).
\let\errorstopmode\nonstopmode
\let\scrollmode\nonstopmode

% ── Font resolution: prefer TeX Live's own copies over OS-installed ones ──────
% luaotfload resolves a font *by name* against a database built from several
% locations, searched in `location_precedence` order – whose default is
% {"system", "texmf", "local"}. On macOS, "system" means ~/Library/Fonts and
% /Library/Fonts – regular Font Book directories any app (or a user, installing
% a font by hand) can add to, independently of TeX Live. If a font of the same
% name ends up there too, it can be a *different, non-identical* copy of what
% TeX Live ships under texmf-dist (different glyph-index ordering internally,
% despite matching filenames). Since reflowtex addresses glyphs by the index
% LuaTeX assigns while compiling but *serves the browser the texmf-dist file*
% (via kpsewhich, see src/pipeline/fonts/fonts.ts), a "system" font winning
% that lookup makes every glyph index wrong – wrong enough that
% normaliseGlyphAddressing (src/pipeline/fonts/addressing.ts) treats it as a
% substituted glyph and PUA-rewrites it, and the
% served font (lacking the PUA codepoints, since they were computed against a
% different file) can't display it at all. Reordering precedence so texmf wins
% keeps compilation and serving on the same file. This can't be set via a
% config file shipped in the repo (luaotfload reads its config before this
% document exists), so it's forced here, before any font is loaded. Only an
% issue if a same-named font happens to exist outside texmf-dist too (nothing
% to do with a stock TeX Live install), but forcing texmf first costs nothing
% when there's no such collision.
\directlua{
  if fonts and fonts.names and fonts.names.set_location_precedence then
    fonts.names.set_location_precedence({'texmf', 'system', 'local'})
    fonts.names.flush_lookup_cache()
  end
}

% The directive above cannot reach one font: the *default* text face,
% \TU/lmr/m/n/10. LuaLaTeX loads it while the format is built (that is
% where "PRELOADED" fonts come from), long before this document exists, so it
% was resolved under the default precedence -- and NFSS keeps reusing that
% definition for every later \selectfont. Nearly every block is set in it.
% Drop the NFSS identifier of each OpenType face loaded so far and re-select,
% so the default face is looked up again under the precedence just set. A
% fresh load of any *other* face was never affected. Backslashes are built
% with string.char so TeX does not tokenize them inside \directlua.
\directlua{
  local bs = string.char(92)
  for id, f in font.each() do
    local ident = tex.fontidentifier(id)
    if ident and ident:sub(1, 4) == bs .. "TU/" then
      texio.write_nl("reflowtex: re-resolving preloaded font " .. ident)
      tex.print(bs .. "expandafter" .. bs .. "let" .. bs .. "csname "
                .. ident:sub(2) .. bs .. "endcsname" .. bs .. "relax")
    end
  end
}
\selectfont

\makeatletter
% ── After the class ────────────────────────────────────────────────────────
% Code for when the document's class is loaded: at the end of the outermost
% class (a class may load another, with \LoadClass, whose end comes first).
\newcount\reflowtex@classdepth
\let\reflowtex@afterclass\@empty
\def\reflowtexAfterClass#1{\g@addto@macro\reflowtex@afterclass{#1}}
\def\reflowtex@runafterclass{%
  \global\let\reflowtex@afterclass@run\reflowtex@afterclass
  \global\let\reflowtex@afterclass\@empty
  \reflowtex@afterclass@run}
\AddToHook{class/before}{\global\advance\reflowtex@classdepth\@ne}
\AddToHook{class/after}{%
  \global\advance\reflowtex@classdepth\m@ne
  \ifnum\reflowtex@classdepth=\z@ \expandafter\reflowtex@runafterclass\fi}
\makeatother

% Reflow TeX produces a continuous, pageless document: a float is set where it
% is written, as float's [H] sets one. Left to LaTeX, a float waits for room on
% a page that never fills: one that may go "here" is put back by the output
% routine, often after the prose that follows it, and one that may not (a [p]
% figure, a type with no h among its placements) waits for the final
% \clearpage and goes out on a page of its own, which is shipped, not
% captured – lost.
%
% Every float is begun by the kernel's \@xfloat, whatever made it: figure and
% table and their starred forms, a class's own, float's \newfloat and
% \restylefloat (algorithm's), newfloat's \DeclareFloatingEnvironment,
% algorithm2e, rotating, listings. \@xfloat builds the float's box as LaTeX
% does, and the class's or package's end of the float finishes it (float's
% rules and caption above, say). Here \@floatpenalty is then 0, which every
% end of a float (the kernel's, revtex's, acmart's) takes for a float it
% cannot queue: none is queued. The box goes into the flow when the
% environment's group closes, between \intextsep above and below; a float
% written in a paragraph ends it. Its register goes back to the free list, as
% the output routine would give it back. Placement options are read, as LaTeX
% reads them, and have no effect.
\makeatletter
% The kernel's (or the class's) \@xfloat and \@floatboxreset are wrapped at
% the end of the class.
\def\reflowtex@newxfloat#1[#2]{%
  \ifinner
    % LaTeX's "not in outer par mode" error: nothing to place
    \expandafter\reflowtex@xfloat
  \else
    \aftergroup\reflowtex@placefloat
    \expandafter\reflowtex@xfloatbox
  \fi
  {#1}[#2]}
\def\reflowtex@xfloatbox#1[#2]{%
  \reflowtex@xfloat{#1}[#2]%
  % inside the float's box now
  \global\@floatpenalty\z@
  \xdef\reflowtex@floatbox{\@currbox}}
\def\reflowtex@placefloat{%
  \par\addvspace{\intextsep}%
  \box\reflowtex@floatbox
  \@cons\@freelist\reflowtex@floatbox
  \addvspace{\intextsep}}
% A float's caption is a paragraph, small and centred, spaced as LaTeX's
% \@makecaption spaces it (\abovecaptionskip above, \belowcaptionskip below,
% whatever the class set them to). LaTeX's own sets a caption shorter than
% the line as a box as wide as it, which the viewer could not break again at
% another width. \caption itself is the kernel's (or the class's, or the
% caption package's): the counter, the list of figures, the label and the
% anchor are as LaTeX makes them. Set in every float's box, [H]'s too;
% \captionof outside a float keeps the class's.
\def\reflowtex@newfloatboxreset{%
  \reflowtex@floatboxreset
  \let\@makecaption\reflowtex@makecaption}
\long\def\reflowtex@makecaption#1#2{%
  \vskip\abovecaptionskip
  {\small\centering#1: #2\par}%
  \vskip\belowcaptionskip}
\reflowtexAfterClass{%
  \let\reflowtex@xfloat\@xfloat \let\@xfloat\reflowtex@newxfloat
  \let\reflowtex@floatboxreset\@floatboxreset \let\@floatboxreset\reflowtex@newfloatboxreset}
\makeatother

% ── reflowtex node-list serializer ─────────────────────────────────────────────
% Exports the finished node list as output.json (see src/extract/serializer.lua).
\directlua{dofile("serializer.lua")}

% ── Protrusion and expansion codes changing ──────────────────────────────────
% TeX breaks each paragraph with the codes its fonts hold at that moment, and
% they may change along the document (microtype sets a font's anew for each
% NFSS shape it serves). The serializer reads a font's codes again only when
% they may have changed, so the primitives that change them say each time
% they are used: each is a macro that counts (Serializer.codes_changed, a Lua
% function called by number, expandable as \directlua is) and is then the
% primitive, which TeX goes on to read or assign as before. See
% serializer.lua, "Protrusion and expansion".
\makeatletter
\directlua{
  local n = luatexbase.new_luafunction("reflowtex.codes_changed")
  lua.get_functions_table()[n] = Serializer.codes_changed
  token.set_macro("reflowtex@codesfunction", tostring(n), "global")}
\def\reflowtex@countuses#1{%
  \expandafter\let\csname reflowtex@primitive@#1\expandafter\endcsname\csname #1\endcsname
  \expandafter\edef\csname #1\endcsname{\noexpand\luafunction\reflowtex@codesfunction\space
    \expandafter\noexpand\csname reflowtex@primitive@#1\endcsname}}
\reflowtex@countuses{lpcode}\reflowtex@countuses{rpcode}\reflowtex@countuses{efcode}%
\reflowtex@countuses{expandglyphsinfont}
\makeatother

% Nothing here ships pages in the normal course: the serializer keeps the page
% from filling. An output routine that never settles against that (a class's
% own page grid) ships empty pages without end instead, so a run is stopped
% at 2000 pages, with an error the pipeline reports, rather than left to run
% for an hour. (Captured pictures ship through the primitive, uncounted. The
% run ends from Lua: \end is not allowed inside an output routine.)
\AddToHook{shipout/after}{%
  \ifnum\ReadonlyShipoutCounter>2000
    \PackageError{reflowtex}{More than 2000 pages shipped: the document's
      output routine does not settle}{}%
    \directlua{os.exit(1)}%
  \fi}

% ── Pictures ──────────────────────────────────────────────────────────────
% Every picture (a TikZ picture captured, an \includegraphics recorded) is
% numbered by this counter.
\newcount\reflowtexPicCount

% Ordinary \includegraphics images arrive as the same kind of LuaTeX image
% rule, and likewise do not expose their filename through the node API.
% graphicx hands every file it includes, with the name it found it under
% (through \graphicspath, \input@path or with a default extension), to
% \Gin@setfile. The call is wrapped here: it records that name and the PDF
% page asked for, and stamps the image it makes with the record's id.
% Installed when graphics is loaded (graphicx loads it), by whatever loads it
% and with whatever options: the document's own, a class's, TikZ's.
\makeatletter
\def\reflowtex@Ginsetfile#1#2#3{%
  \global\advance\reflowtexPicCount by 1
  \directlua{Serializer.note_graphic(\the\reflowtexPicCount,
    "\luaescapestring{#3}", "\luaescapestring{\Gin@page}")}%
  \begingroup\attribute902=\reflowtexPicCount\relax
    \reflowtexOrigGinSetfile{#1}{#2}{#3}%
  \endgroup}
\AddToHook{package/graphics/after}{%
  \let\reflowtexOrigGinSetfile\Gin@setfile
  \let\Gin@setfile\reflowtex@Ginsetfile}
\makeatother

% ── Companion package switch ─────────────────────────────────────────────────
% \ifreflowtex is true only inside this extraction run. reflowtex.sty, which a
% snippet may load, defines the same switch as *false* when it finds it
% undefined, so one source compiles to a PDF unchanged and to a web build with
% the web-only constructs live. Defined before the preamble so the package,
% loaded from the preamble, sees it.
\newif\ifreflowtex
\reflowtextrue

% Streams (\begin{webstream} in reflowtex.sty, and batch parts below)
% are numbered by this one counter, which the package uses when it finds it.
\newcount\reflowtexStreamCount

% ── Asides ──────────────────────────────────────────────────────────────────
% \reflowtexAside{kind}{key=value,…}{text}: `text` typeset where it stands, in
% running text too, but kept out of the flow – a stream of `kind`, which the
% serializer marks aside=true and the viewer shows on demand (a margin note,
% a popover). Like a footnote's text it goes into an insertion (attribute 914
% on the insertion, 911 on its contents); in running text it also leaves an
% empty box where it stood (attribute 915), which flows with the text, for
% the viewer to place the aside by. \marginpar (reflowtex-capture-late.tex) and reflowtex.sty's
% \webaside are made of it. Parameters are taken as written: expand them first.
% Its text is set as LaTeX sets a \marginpar or a footnote, with the paragraph
% settings restored (\@parboxrestore): what holds around it – a list's margins,
% \raggedright's \\ – does not apply inside.
\newinsert\reflowtexAsideIns
\count\reflowtexAsideIns=0 \dimen\reflowtexAsideIns=\maxdimen \skip\reflowtexAsideIns=0pt
\makeatletter
\long\def\reflowtexAside#1#2#3{%
  \global\advance\reflowtexStreamCount by 1
  \directlua{Serializer.note_stream(\the\reflowtexStreamCount,
    "\luaescapestring{#1}", 0, "\luaescapestring{#2}")}%
  \begingroup
    \ifhmode\attribute915=\reflowtexStreamCount\relax\hbox{}\attribute915=-"7FFFFFFF\fi
    \attribute914=\reflowtexStreamCount\relax
    \insert\reflowtexAsideIns{%
      \attribute911=\reflowtexStreamCount\relax
      \hsize\linewidth \@parboxrestore #3\par}%
  \endgroup}
\makeatother

% ── Batch parts ──────────────────────────────────────────────────────────────
% Pipeline.compile_batch compiles several snippets (a book's chapters) as one
% document, joined by \reflowtexbatchpart{N}. Everything typeset after it
% belongs to stream "batch-part" N, by a *global* attribute assignment rather
% than an environment, so no group is opened: a \newcommand in one part is
% still defined in the next, as in the book. The pipeline cuts the finished
% document at these streams, one bundle per part.
\def\reflowtexbatchpart#1{%
  \par
  \global\advance\reflowtexStreamCount by 1
  \directlua{Serializer.note_stream(\the\reflowtexStreamCount, "batch-part", -1, "part=#1")}%
  \global\attribute911=\reflowtexStreamCount\relax}

% The width increment the driver sets (display sampling, the render tests);
% applied at \begin{document} (reflowtex-capture-late.tex).
\newdimen\reflowtexWidthExtra

% Capture the finished TikZ hbox rather than scanning its source. This also
% handles front ends such as tikz-cd, which call the low-level \tikzpicture and
% \endtikzpicture macros and are incompatible with PGF's LaTeX source scanner.
% A copy is shipped as a tightly-sized private PDF page; the document flow gets
% an empty box with identical metrics and attribute 908. serializer.lua turns
% that placeholder into a picture node referring to the private page.
%
% If the document enabled PGF externalisation, it is turned off first: capturing boxes
% is the single picture path, independent of how the source was configured.
%
% Only the outermost picture is captured. A picture inside another – a \tikz in
% a node's text, or the one pgfplots' axis opens inside the user's – is part of
% the outer box, and its drawing and labels are captured with it. Captured on
% its own, it left an empty placeholder in the outer box, which is shipped, not
% serialized: the inner drawing was lost, and a pgfplots axis came out blank.
\makeatletter
% The machinery is defined here; \tikzpicture is wrapped at \begin{document}
% (reflowtex-capture-late.tex), when the preamble has loaded TikZ, if it has.
\newbox\reflowtexTikzBox
\newbox\reflowtexTikzPageBox
\newdimen\reflowtexTikzWidth
\newdimen\reflowtexTikzHeight
\newdimen\reflowtexTikzDepth
\newcount\reflowtexTikzNesting
\ExplSyntaxOn
\cs_new_eq:NN \reflowtexPrimitiveShipout \tex_shipout:D
\ExplSyntaxOff
\def\reflowtex@tikzpicture{%
  \global\advance\reflowtexTikzNesting by 1
  \ifnum\reflowtexTikzNesting=1
    \expandafter\reflowtexCaptureTikzpicture
  \else
    \expandafter\reflowtexOrigTikzpicture
  \fi}
\def\reflowtex@endtikzpicture{%
  \global\advance\reflowtexTikzNesting by -1
  \ifnum\reflowtexTikzNesting=0
    \expandafter\reflowtexCaptureEndtikzpicture
  \else
    \expandafter\reflowtexOrigEndtikzpicture
  \fi}
\def\reflowtexCaptureTikzpicture{\reflowtexCaptureBegin\reflowtexOrigTikzpicture}
\def\reflowtexCaptureEndtikzpicture{\reflowtexOrigEndtikzpicture\reflowtexCaptureEnd}
% The capture itself, for anything PGF draws: \reflowtexCaptureBegin opens the
% box, \reflowtexCaptureEnd closes it, ships it and leaves the placeholder. In
% maths the box is a formula in the style it was begun in (LuaTeX's
% \mathstyle), so a matrix keeps its size and its axis.
\newif\ifreflowtexCaptureMath
\def\reflowtexCaptureBegin{%
  \global\advance\reflowtexPicCount by 1
  \edef\reflowtexThisPic{\the\reflowtexPicCount}%
  \ifmmode
    \reflowtexCaptureMathtrue
    \edef\reflowtexCaptureStyle{\ifcase\numexpr\mathstyle/2\relax
      \noexpand\displaystyle\or\noexpand\textstyle\or\noexpand\scriptstyle\else\noexpand\scriptscriptstyle\fi}%
    \setbox\reflowtexTikzBox=\hbox\bgroup$\reflowtexCaptureStyle\m@th
  \else
    \reflowtexCaptureMathfalse
    \setbox\reflowtexTikzBox=\hbox\bgroup
  \fi}
\def\reflowtexCaptureEnd{%
  \ifreflowtexCaptureMath$\fi
  \egroup
  \reflowtexTikzWidth=\wd\reflowtexTikzBox
  \reflowtexTikzHeight=\ht\reflowtexTikzBox
  \reflowtexTikzDepth=\dp\reflowtexTikzBox
  \begingroup
    \pagewidth=\reflowtexTikzWidth
    \pageheight=\dimexpr\reflowtexTikzHeight+\reflowtexTikzDepth\relax
    % the page is the box, wherever the document's pages are moved on the
    % paper: by \hoffset and \voffset (a class's \advance\hoffset-1in,
    % geometry's hoffset) or by the PDF origin
    \hoffset=0pt
    \voffset=0pt
    \pdfvariable horigin=0pt
    \pdfvariable vorigin=0pt
    \setbox\reflowtexTikzPageBox=\vbox{%
      \offinterlineskip
      \hbox{\copy\reflowtexTikzBox}}%
    % shipping runs the \writes inside the picture (an \addcontentsline in a
    % node or a box): as LaTeX's own shipout does, \protect leaves what
    % follows it unexpanded
    \let\protect\noexpand
    \reflowtexPrimitiveShipout\box\reflowtexTikzPageBox
    \directlua{Serializer.note_captured_picture(
      \reflowtexThisPic, status.total_pages)}%
  \endgroup
  % PGF places a finished picture with \leavevmode (\pgfsys@typesetpicturebox),
  % so a picture that opens a paragraph – alone, or in a center – starts the
  % paragraph. Here that \leavevmode ran inside the capture box above, so the
  % placeholder has to do it itself, or it would join the vertical list
  % outside any paragraph and never reach the page.
  \leavevmode
  \begingroup
    \attribute908=\reflowtexThisPic\relax
    \setbox\reflowtexTikzBox=\hbox to \reflowtexTikzWidth{\hss}%
    \ht\reflowtexTikzBox=\reflowtexTikzHeight
    \dp\reflowtexTikzBox=\reflowtexTikzDepth
    \box\reflowtexTikzBox
  \endgroup}

% nicematrix draws with PGF too, but not in one picture: every cell has a
% small picture of its own, remembering its place on the page (remember
% picture), and the rules (hvlines), \Block contents, dotted lines and
% \CodeAfter's drawing are pictures placed from those positions. Captured
% cell by cell, each would be on a page of its own and nothing could find its
% place; left in the text, the drawing (PDF literals) is lost. So the whole
% environment is the picture, captured like the outermost TikZ picture (a
% TikZ picture inside it is part of it). Its positions come from the .aux,
% as in any document: the run is repeated until they settle (lualatex.ts).
\def\reflowtexCaptureEnvBegin{%
  \global\advance\reflowtexTikzNesting by 1
  \ifnum\reflowtexTikzNesting=1 \expandafter\reflowtexCaptureBegin\fi}
\def\reflowtexCaptureEnvEnd{%
  \global\advance\reflowtexTikzNesting by -1
  \ifnum\reflowtexTikzNesting=0 \expandafter\reflowtexCaptureEnd\fi}
\def\reflowtexCaptureEnvironment#1{%
  \AddToHook{env/#1/before}{\reflowtexCaptureEnvBegin}%
  \AddToHook{env/#1/after}{\reflowtexCaptureEnvEnd}}
\@for\reflowtex@env:={NiceTabular,NiceTabular*,NiceTabularX,NiceArray,NiceArrayWithDelims,%
    NiceMatrix,pNiceMatrix,bNiceMatrix,BNiceMatrix,vNiceMatrix,VNiceMatrix,%
    pNiceArray,bNiceArray,BNiceArray,vNiceArray,VNiceArray}\do
  {\expandafter\reflowtexCaptureEnvironment\expandafter{\reflowtex@env}}
\makeatother
